\chapter*{Anexo A — Código fuente del simulador Arduino}
\addcontentsline{toc}{chapter}{Anexo A — Código fuente del simulador Arduino}
\label{anexo:arduino}

Este anexo recoge los fragmentos más representativos del código fuente del simulador
de ECU implementado sobre Arduino MKR. El código completo está disponible en el
repositorio del proyecto.

\section*{A.1 — Definiciones de protocolo (\texttt{ecu\_protocol.h})}

El fichero de cabecera centraliza todas las constantes del protocolo, las estructuras
de datos y las funciones de codificación/decodificación OBD-II.

\begin{lstlisting}[language=C++, caption={Constantes de CAN, ISO-TP y ruido — ecu\_protocol.h}]
// ---------- CAN BUS ----------
#define CAN_SPEED          500E3
#define ECU_CAN_ID         0x7E0
#define ECU_RESPONSE_ID    0x7E8

// ---------- ISO-TP ----------
#define ISOTP_PCI_SF       0x00  // Single Frame
#define ISOTP_PCI_FF       0x10  // First Frame
#define ISOTP_PCI_CF       0x20  // Consecutive Frame
#define ISOTP_PCI_FC       0x30  // Flow Control

#define ISOTP_CF_SEP_MS      25
#define ISOTP_SF_MAX_PAYLOAD 7
#define ISOTP_FF_DATA_BYTES  6

// ---------- CAN NOISE ----------
#define CAN_NOISE_DROP_PCT         2
#define CAN_NOISE_NRC_PCT          1
#define CAN_NOISE_LATENCY_MIN_MS   1
#define CAN_NOISE_LATENCY_MAX_MS   15
#define CAN_NOISE_BG_TRAFFIC_PCT   30
static const uint32_t CAN_BG_IDS[4] = {0x280, 0x480, 0x320, 0x520};

// ---------- HARDWARE ----------
#define DTC_FAULT_PIN        6
#define DTC_DEBOUNCE_MS      300
\end{lstlisting}

\begin{lstlisting}[language=C++, caption={Estructura VehicleData y funciones de codificación OBD-II}]
struct VehicleData {
  char     vin[18];
  uint16_t rpm;
  uint8_t  speed;
  int16_t  coolantTemp;
  uint8_t  throttlePos;
  uint8_t  engineLoad;
  uint16_t batteryVoltage;  // mV
  bool     checkEngine;
  uint8_t  numDTCs;
};

// Codificacion OBD-II: rpm*4 en dos bytes big-endian
inline void encodeRPM(uint8_t* data, uint16_t rpm) {
  uint16_t encoded = rpm * 4;
  data[0] = (encoded >> 8) & 0xFF;
  data[1] = encoded & 0xFF;
}

// Temperatura: T_raw = T_celsius + 40
inline uint8_t encodeTemp(int16_t temp) {
  return (uint8_t)(temp + 40);
}

// Porcentaje: val_raw = percent * 255 / 100
inline uint8_t encodePercent(uint8_t percent) {
  return (percent * 255) / 100;
}
\end{lstlisting}

\section*{A.2 — Bucle principal y mecanismos de ruido CAN (\texttt{ecu\_sim.ino})}

\begin{lstlisting}[language=C++, caption={Interrupción hardware para inyección de DTCs e inicialización}]
volatile bool dtcFaultPressed  = false;
unsigned long lastDtcDebounce  = 0;
bool canNoiseEnabled = true;

void onDtcFaultButton() {
  dtcFaultPressed = true;
}

void setup() {
  Serial.begin(SERIAL_BAUDRATE);
  CAN.begin(CAN_SPEED);

  pinMode(DTC_FAULT_PIN, INPUT_PULLUP);
  attachInterrupt(digitalPinToInterrupt(DTC_FAULT_PIN),
                  onDtcFaultButton, FALLING);

  initVehicleData();
}

void loop() {
  updatePhysicsModel();
  handleDtcFaultButton();
  generateBackgroundTraffic();

  int packetSize = CAN.parsePacket();
  if (packetSize > 0) {
    processCANMessage();
  }
}
\end{lstlisting}

\begin{lstlisting}[language=C++, caption={Inyección de DTC por interrupción hardware con debounce}]
static const uint16_t DTC_POOL[] = {
  0x0300, 0x0301, 0x0302, 0x0303,
  0x0113, 0x0118, 0x0340, 0x0341,
  0x0500, 0x0501, 0x0171, 0x0172
};
#define DTC_POOL_SIZE 12

void handleDtcFaultButton() {
  if (!dtcFaultPressed) return;
  dtcFaultPressed = false;

  unsigned long now = millis();
  if (now - lastDtcDebounce < DTC_DEBOUNCE_MS) return;
  lastDtcDebounce = now;

  uint16_t candidate = DTC_POOL[random(DTC_POOL_SIZE)];

  // Deduplicacion: no insertar si ya esta presente
  for (uint8_t i = 0; i < numStoredDTCs; i++) {
    if (dtcList[i] == candidate) return;
  }

  if (numStoredDTCs < MAX_DTCS) {
    dtcList[numStoredDTCs++] = candidate;
    checkEngineLight = true;
  }
}
\end{lstlisting}

\begin{lstlisting}[language=C++, caption={Mecanismos de ruido CAN dentro del procesador de mensajes}]
void processCANMessage() {
  uint8_t buf[8];
  int len = 0;
  while (CAN.available() && len < 8) {
    buf[len++] = CAN.read();
  }

  if (canNoiseEnabled) {
    // Latencia variable: simula carga del bus
    delay(random(CAN_NOISE_LATENCY_MIN_MS, CAN_NOISE_LATENCY_MAX_MS));

    // Descarte aleatorio de paquetes (2%)
    if (random(100) < CAN_NOISE_DROP_PCT) return;

    // Respuesta negativa UDS esporadica (1%)
    if (random(100) < CAN_NOISE_NRC_PCT) {
      sendNegativeResponse(buf[1], NRC_CONDITIONS_NOT_CORRECT);
      return;
    }
  }

  // Decodificacion del modo y PID
  uint8_t mode = buf[1];
  uint8_t pid  = buf[2];
  dispatchMode(mode, pid);
}

void generateBackgroundTraffic() {
  if (!canNoiseEnabled) return;
  if (random(100) >= CAN_NOISE_BG_TRAFFIC_PCT) return;

  uint32_t id = CAN_BG_IDS[random(CAN_BG_ID_COUNT)];
  CAN.beginPacket(id);
  for (int i = 0; i < 8; i++) CAN.write(random(256));
  CAN.endPacket();
}
\end{lstlisting}

\begin{lstlisting}[language=C++, caption={Respuesta VIN con secuencia ISO-TP multi-frame}]
void sendVIN() {
  const char* vin = vehicleData.vin;  // 17 bytes ASCII

  // First Frame: PCI 0x10 0x14 (longitud total = 20 bytes)
  uint8_t ff[8] = {
    0x10, 0x14,          // FF PCI: len=20
    0x49, 0x02, 0x01,    // modo respuesta, PID, count
    (uint8_t)vin[0], (uint8_t)vin[1], (uint8_t)vin[2]
  };
  CAN.beginExtendedPacket(ECU_RESPONSE_ID);
  for (int i = 0; i < 8; i++) CAN.write(ff[i]);
  CAN.endPacket();

  // Espera Flow Control del tester
  waitFlowControl();

  // Consecutive Frame 1: bytes VIN[3..9]
  uint8_t cf1[8] = {0x21};
  for (int i = 0; i < 7; i++) cf1[i+1] = (uint8_t)vin[3+i];
  delay(ISOTP_CF_SEP_MS);
  CAN.beginExtendedPacket(ECU_RESPONSE_ID);
  for (int i = 0; i < 8; i++) CAN.write(cf1[i]);
  CAN.endPacket();

  // Consecutive Frame 2: bytes VIN[10..16] + padding
  uint8_t cf2[8] = {0x22};
  for (int i = 0; i < 7; i++) {
    cf2[i+1] = (10+i < VIN_LENGTH) ? (uint8_t)vin[10+i] : ISOTP_PADDING_BYTE;
  }
  delay(ISOTP_CF_SEP_MS);
  CAN.beginExtendedPacket(ECU_RESPONSE_ID);
  for (int i = 0; i < 8; i++) CAN.write(cf2[i]);
  CAN.endPacket();
}
\end{lstlisting}
